%%% SECTION 1: INTRODUCTION TO PHYSICAL AI ONCOLOGY TRIALS %%%

\subsection{Background and Motivation}
\label{subsec:background}

Oncology drug development faces severe challenges in cost, timeline, and
patient access. The Tufts Center for the Study of Drug Development
\cite{tufts-delay} has quantified the value of each day of delay in drug
development, demonstrating that even modest reductions in development
timelines translate to significant economic and patient benefit. Current
oncology clinical trials, as tracked through ClinicalTrials.gov
\cite{clinicaltrials-gov} and described by the National Cancer Institute
\cite{nci-clinical-trials}, operate within a framework that was designed for
human-directed operations at every stage. The U.S. Food and Drug
Administration provides extensive guidance on clinical trial conduct
\cite{fda-clinical-trials-guidance} and oncology-specific considerations
\cite{fda-oncology-guidance}, yet these resources address AI and robotics
only through individual guidance documents that treat each technology in
isolation.

The fundamental problem is structural: oncology trials require coordination
across multiple regulatory frameworks (IND applications under 21 CFR Part 312
\cite{cfr-part-312}, human subjects protection under 21 CFR Part 50
\cite{cfr-part-50}, good clinical practice under ICH E6(R3) \cite{ich-e6r3}),
multiple stakeholders (sponsors, investigators, IRBs, sites, patients), and
multiple technology systems (electronic data capture, imaging, laboratory
information systems). Each additional technology layer increases complexity.
Physical AI introduces robotic systems that perform functions traditionally
reserved for human clinical staff, creating coordination challenges that exceed
the capacity of piecemeal regulatory guidance.

The constitutional framework of U.S. government, spanning legislative (Article
I), executive (Article II), and judicial (Article III) branches, provides the
institutional infrastructure within which clinical trials operate. The
Administrative Procedure Act (5 U.S.C. \S~551 et seq.) governs FDA rulemaking.
The Freedom of Information Act (5 U.S.C. \S~552) ensures transparency. The
Privacy Act (5 U.S.C. \S~552a) protects individual data in federal systems.
HIPAA \cite{hipaa} governs health information privacy and security across
covered entities and business associates. These statutes were not designed
with Physical AI in mind, but they provide a robust foundation upon which
Physical AI regulation can be built through targeted adaptation rather than
wholesale replacement.

The California and federal regulatory landscape presents a dual-jurisdiction
challenge for Physical AI trial sites. California's Protection of Human
Subjects in Medical Experimentation Act (Cal. Health \& Safety Code
\S\S~24170-24179.5), the Confidentiality of Medical Information Act (CMIA,
Cal. Civ. Code \S\S~56-56.37), and consumer privacy protections under
CCPA/CPRA create additional obligations beyond federal requirements. Federal
requirements include 21 CFR Parts 50, 56, 312, and 812 for trial conduct,
21 CFR Part 11 for electronic records, and 45 CFR Parts 46, 160, and 164 for
human subjects and health information protections. Navigating this dual
framework requires a comprehensive approach that no existing guidance document
provides.

The evidence supporting Physical AI in oncology is substantial. The 24-hour
clinical trial simulation from the Clinical Trial Site Complete Documentation
\cite{site-docs} demonstrated that 168 patients could be treated with 29
robots across 15 cancer types at 99.7 percent uptime with zero patient harm
events. This simulation established that more patients can be treated with
more robots and fewer workers while maintaining safety standards documented
per ICH E6(R3) requirements.

\subsection{Defining Physical AI in Oncology}
\label{subsec:defining-pai}

Physical AI is defined as the integration of embodied robotic systems with
advanced artificial intelligence for direct patient care in clinical trial
settings. This definition encompasses robotic systems that physically interact
with patients, the AI algorithms that guide robotic operations, and the
infrastructure connecting these systems to clinical trial management platforms.
The adapted 21 CFR Part 312 \cite{cfr312-adapt} provides 22 Physical AI-specific
definitions in \S~312.3, establishing standardized terminology for the entire
field. These definitions include Physical AI system, Physical AI adverse event,
Physical AI operator, robotic procedure, autonomous mode, supervised mode, and
Physical AI safety matrix, among others. The adapted 21 CFR Part 50
\cite{cfr50-adapt} provides 17 additional definitions in Subpart A, addressing
terms specific to human subjects protection in robotic contexts.

The FDA Digital Health and Artificial Intelligence Glossary
\cite{fda-ai-glossary} provides official FDA terminology for AI and digital
health technologies, but does not address the specific integration of embodied
robotic systems with AI in clinical trial settings. This gap in official
terminology is one reason why the adapted regulatory standards in this
platform provide their own comprehensive definition sets. Physical AI extends
beyond software-based AI (such as clinical decision support or medical image
analysis) by incorporating physical embodiment: robots that move through
clinical spaces, interact with patients through touch and proximity, and
perform procedures that affect patient anatomy.

State-of-the-art AI systems that enable Physical AI development include Claude
Opus 4.6 \cite{claude-opus}, GPT-5.4 \cite{gpt-5-4}, and Gemini
\cite{gemini}. These large language models serve multiple roles in Physical AI
oncology trials: generating and reviewing clinical documentation, planning
robotic task sequences, performing triple AI peer review of system updates, and
enabling agentic AI workflows where AI systems coordinate multi-step clinical
operations autonomously. The federated learning framework \cite{fl-paper}
leverages these AI systems for its triple peer review pipeline, ensuring that
model updates are validated by multiple independent AI systems before
deployment across trial sites.

Physical AI in oncology spans ten robot categories as documented in Patient
Instructions: Physical AI Trials \cite{patient-instructions-paper}: surgical
robots (da Vinci, Hugo, Versius), collaborative robots (Franka Panda, Kinova
Gen3, UFactory xArm7), radiotherapy positioning robots, needle-placement
robots, social companion robots, humanoid robots (Tesla Optimus, Digit, Boston
Dynamics Atlas), radiotherapy motion-tracking robots, imaging robots, steerable
needle robots, and rehabilitation exoskeletons. Each category serves a distinct
clinical function and is evaluated through the Unification Standard Level (USL)
framework \cite{usl-standard} across four dimensions of readiness.

\subsection{The Case for a National Platform}
\label{subsec:case-national}

This National Platform is more comprehensive and meaningful compared to
existing FDA approaches for several structural reasons. The FDA's approach to
AI in clinical trials operates through individual guidance documents: the
January 2025 draft guidance on AI for drug and biologic regulatory
decision-making proposes a risk-based credibility assessment framework; the
January 2025 draft on AI-enabled device software functions addresses lifecycle
management; the September 2024 final guidance on decentralized trial elements
covers remote monitoring; and the August 2025 final guidance on Predetermined
Change Control Plans (PCCPs) addresses iterative AI device improvement. While
each document is valuable individually, none addresses the integrated
challenge of deploying Physical AI across an entire clinical trial lifecycle.

This platform provides five categories of capability that no existing guidance
matches: (a) three simultaneously adapted core regulatory standards (ICH
E6(R3), 21 CFR Part 50, 21 CFR Part 312) that build on existing
well-established frameworks and add credibility to Physical AI implementations
through the Adaption: ICH Harmonised Guideline \cite{ich-adapt}, Physical AI
Adaption: 21 CFR Part 50 \cite{cfr50-adapt}, and Physical AI Adaption: 21 CFR
Part 312 \cite{cfr312-adapt}; (b) complete site establishment documentation
with eleven regulatory documents covering state legislation, municipal codes,
federal compliance, building standards, operational procedures, and emergency
preparedness \cite{site-docs}; (c) validated simulation evidence at both
single-patient scale (ten-stage pipeline \cite{patient-journey-paper}) and
168-patient scale (24-hour continuous operations \cite{site-docs});
(d) quantitative robot readiness standards through PSL (three dimensions for
site compliance) and USL (four dimensions for robot evaluation) with scores
for nine robots across three categories \cite{usl-standard}; and (e) national
infrastructure for multi-site coordination through five MCP servers
\cite{national-mcp-paper} and federated learning \cite{fl-paper}.

The fragmentation problem in current clinical trial systems, as detailed in the
Physical AI National MCP Servers paper \cite{national-mcp-paper}, illustrates
why an integrated platform is necessary. Current systems operate in silos with
incompatible data formats, disconnected communication channels, and no
standardized protocol for real-time inter-system coordination. Physical AI
magnifies this fragmentation because robots, AI systems, clinical databases,
and regulatory reporting tools must all communicate seamlessly. The five-server
MCP architecture provides standardized communication pathways, while federated
learning enables multi-site model improvement without centralizing patient data.

Even established institutions recognize the need for modernization. The
National Cancer Institute's modernization efforts \cite{nci-modernizing} and
FDA's adaptive designs guidance \cite{fda-adaptive-designs} acknowledge that
current trial infrastructure must evolve. However, none has proposed a
solution as comprehensive as this National Platform, which addresses
regulatory, operational, technical, and economic dimensions simultaneously.

\subsection{Scope and Organization of This Document}
\label{subsec:scope}

This document comprises sixteen sections organized into a logical progression
from regulatory foundations through evidence, infrastructure, economics, and
implementation strategy.

Sections 1 through 3 establish the regulatory context. Section 1 (this
section) introduces Physical AI oncology trials and makes the case for a
national platform. Section 2 presents the U.S. government framework spanning
all three branches, adapted from Research A on federal branch operations.
Section 3 analyzes the dual California and federal regulatory landscape,
adapted from Research B on regulatory requirements.

Sections 4 through 6 present the three adapted regulatory standards. Section 4
covers the Adaption: ICH Harmonised Guideline \cite{ich-adapt} for good
clinical practice. Section 5 covers the Physical AI Adaption: 21 CFR Part 50
\cite{cfr50-adapt} for human subjects protection. Section 6 covers the
Physical AI Adaption: 21 CFR Part 312 \cite{cfr312-adapt} for investigational
new drug applications.

Section 7 presents the quantitative standards frameworks (PSL and USL) that
evaluate site compliance and robot readiness respectively. Section 8 details
the complete documentation package for clinical trial site establishment,
drawing from eleven regulatory documents \cite{site-docs}.

Sections 9 and 10 present the patient-centered evidence. Section 9 covers A
Cancer Patient's Journey, Physical AI \cite{patient-journey-paper} with its
ten-stage pipeline. Section 10 covers Patient Instructions: Physical AI Trials
\cite{patient-instructions-paper} for ten robot types.

Sections 11 and 12 present the national infrastructure. Section 11 covers the
five-server MCP architecture \cite{national-mcp-paper}. Section 12 covers
federated learning for multi-site privacy-preserving coordination
\cite{fl-paper}.

Sections 13 and 14 address economics and implementation. Section 13 provides
the financial and economic impact analysis. Section 14 presents the national
implementation strategy with phased deployment.

Sections 15 and 16 synthesize and conclude. Section 15 discusses the
implications, limitations, and significance of the platform. Section 16
presents key findings and a call to action. Five appendices provide reference
material including the complete source file directory, glossary, regulatory
cross-reference matrix, scoring reference tables, and simulation evidence
summaries.

\subsection{Contributions}
\label{subsec:contributions}

The key contributions of this National Platform are:

\begin{enumerate}[nosep]
  \item The first comprehensive end-to-end regulatory framework for Physical AI
    oncology trials, integrating regulatory adaptation, site establishment,
    patient care, and national infrastructure into a unified platform.
  \item Three simultaneously adapted core regulatory standards (ICH E6(R3),
    21 CFR Part 50, 21 CFR Part 312) that build on internationally recognized
    frameworks and add credibility through established practice.
  \item Quantitative PSL (three dimensions) and USL (four dimensions) scoring
    frameworks for objective, reproducible assessment of site compliance and
    robot readiness with scores for nine robots across three categories.
  \item Validated simulation evidence at scale: a single-patient ten-stage
    pipeline and a 24-hour simulation treating 168 patients with 29 robots at
    99.7 percent uptime and zero patient harm events.
  \item Complete site establishment documentation across eleven regulatory
    documents covering legislation, municipal codes, regulations, compliance
    guides, building standards, and operational procedures.
  \item National MCP server architecture with five servers, 23 tools,
    standardized conformance levels, and emergency safety modules for
    nationwide Physical AI trial connectivity.
  \item Federated learning framework with five foundational pillars, triple AI
    peer review, and privacy-preserving multi-site model improvement under
    HIPAA and state privacy laws.
  \item Patient-centered trial design with comprehensive rights, ten-robot
    patient instructions, ongoing consent processes, and demonstrated patient
    benefits in convenience, cost, and quality of care.
  \item Economic analysis demonstrating per-patient cost savings, timeline
    compression, and favorable scale economics from single-site to nationwide
    deployment.
  \item Open-source implementation across three repositories
    \cite{main-repo, mcp-repo, fl-repo} under the MIT license, enabling
    community evaluation, adoption, and extension.
\end{enumerate}
